\documentclass[a4paper,10pt]{ctexart}

\usepackage[margin=2.54cm]{geometry}
\usepackage{newtxmath}
\usepackage{fontspec}
\usepackage{amsmath}
\usepackage{bm}
\usepackage{graphicx}
\usepackage{svg}
\usepackage{fancyhdr}
\usepackage[backend=biber]{biblatex}

% Western text/numerals use Times New Roman and mathematics uses a
% Times-compatible font. Chinese text and headings follow ctexart defaults.
\setmainfont{Times New Roman}
\addbibresource{ref.bib}

\setlength{\parindent}{2em}
\setlength{\parskip}{0pt}
\pagestyle{fancy}
\fancyhf{}
\fancyhead[R]{Continuous Wave Second Harmonics Generation}
\fancyfoot[C]{\thepage}
\renewcommand{\headrulewidth}{0pt}

\newcommand{\dd}{\mathop{}\!\mathrm{d}}

\title{Continuous Wave Second Harmonics Generation}
\author{Wei}
\date{}

\begin{document}

\maketitle

\section{Theory}

\subsection{Coupled wave equations}
Considering diffraction, walk off and nonlinearity, the coupled wave equations are:
\begin{equation}
  \begin{aligned}
    \frac{\partial A_{\mathrm{SH}}}{\partial z}
             & =\frac{\mathrm{i}}{2k_{\mathrm{SH}}}\nabla_\perp^2A_{\mathrm{SH}}
    +\tan\rho\frac{\partial A_{\mathrm{SH}}}{\partial y}
    +\frac{\mathrm{i}\omega_{\mathrm{SH}}}{2n_{\mathrm{SH}}c}
    d_{\mathrm{eff}}A_{\mathrm{F}}^2\exp(-\mathrm{i}\Delta kz),                  \\
    \frac{\partial A_{\mathrm{F}}}{\partial z}
             & =\frac{\mathrm{i}}{2k_{\mathrm{F}}}\nabla_\perp^2A_{\mathrm{F}}
    +\frac{\mathrm{i}\omega_{\mathrm{F}}}{n_{\mathrm{F}}c}
    d_{\mathrm{eff}}A_{\mathrm{SH}}A_{\mathrm{F}}^*
    \exp(\mathrm{i}\Delta kz),                                                   \\
    \Delta k & =k_{\mathrm{SH}}-2k_{\mathrm{F}}.
  \end{aligned}
  \label{eq:CWE}
\end{equation}
For LBO in the XY plane, $d_{\mathrm{eff}}=d_{32}\cos\phi$, where
$d_{32}=-0.98\,\mathrm{pm/V}$.

The input Gaussian beam is:
\begin{equation}
  A(x,y,z)=A_0\frac{\omega_0}{\omega(z)}
  \exp\left[-\frac{x^2+y^2}{\omega^2(z)}\right]
  \exp\left\{\mathrm{i}k\left[z+\frac{x^2+y^2}{2R(z)}\right]
  -\mathrm{i}\arctan\left(\frac{z}{z_R}\right)\right\}.
  \label{eq:gaussian-beam}
\end{equation}
where $\omega_0$ and $\omega(z)$ are beam radii at $1/e^2$ intensity and $z$ is the
distance from the beam waist. If the average power is $P$ and the refractive
index of the medium is $n$, the other parameters are:
\begin{equation}
  \begin{aligned}
    k         & =\frac{2\pi}{\lambda},                           \\
    \omega(z) & =\omega_0\sqrt{1+\left(\frac{z}{z_R}\right)^2},  \\
    R(z)      & =z+\frac{z_R^2}{z},                              \\
    z_R       & =\frac{\pi\omega_0^2}{\lambda},                  \\
    A_0       & =\sqrt{\frac{4P}{\pi\omega_0^2nc\varepsilon_0}}.
  \end{aligned}
\end{equation}

\begin{figure}[t]
  \centering
  \IfFileExists{../fig/beam.svg}{%
    \includesvg[width=0.75\linewidth]{../fig/beam}%
  }{%
    \fbox{\parbox[c][3cm][c]{0.7\linewidth}{\centering
        Missing figure: \texttt{../fig/beam.svg}}}%
  }
  \caption{Gaussian beam is focused to the nonlinear crystal.}
  \label{fig:gaussian-beam}
\end{figure}

The definition of the D4$\sigma$ radius is:
\begin{equation}
  \omega_x=2\sqrt{\frac{\iint x^2I(x,y)\dd x\dd y}
    {\iint I(x,y)\dd x\dd y}}.
\end{equation}

\subsection{Numerical algorithm}
Equation~\eqref{eq:CWE} is solved by using the RK4IP method
\cite{balac_embedded_2013}, but without adaptive step-size control.
Equation~\eqref{eq:CWE} can be organized as follows:
\begin{equation}
  \begin{aligned}
    \frac{\partial A_S}{\partial z}
     & =\mathcal{D}_SA_S+\mathcal{N}_S(A_S,A_F), \\
    \frac{\partial A_F}{\partial z}
     & =\mathcal{D}_FA_S+\mathcal{N}_F(A_S,A_F).
  \end{aligned}
\end{equation}
where $\mathcal{D}$ represents linear effects and $\mathcal{N}$ represents
nonlinear coupling. Then transform $A_S$ and $A_F$ to the interaction picture:
\begin{equation}
  \begin{aligned}
    A_S^{\mathrm{ip}} & =\exp[-(z-z')\mathcal{D}_S]A_S, \\
    A_F^{\mathrm{ip}} & =\exp[-(z-z')\mathcal{D}_F]A_F.
  \end{aligned}
\end{equation}
The coupled equations in the interaction picture are:
\begin{equation}
  \begin{aligned}
    \frac{\partial A_S^{\mathrm{ip}}}{\partial z}
     & =\exp[-(z-z')\mathcal{D}_S]\mathcal{N}_S
    \left(\exp[(z-z')\mathcal{D}_S]A_S^{\mathrm{ip}},
    \exp[(z-z')\mathcal{D}_F]A_F^{\mathrm{ip}}\right), \\
    \frac{\partial A_F^{\mathrm{ip}}}{\partial z}
     & =\exp[-(z-z')\mathcal{D}_F]\mathcal{N}_F
    \left(\exp[(z-z')\mathcal{D}_S]A_S^{\mathrm{ip}},
    \exp[(z-z')\mathcal{D}_F]A_F^{\mathrm{ip}}\right).
  \end{aligned}
  \label{eq:IP-CWE}
\end{equation}
By choosing $z'=z+h/2$, Equation~\eqref{eq:IP-CWE} can be solved by the RK4
method:
\begin{align}
  \alpha_{1S} & =\exp\left(\frac{h}{2}\mathcal{D}_S\right)
  \mathcal{N}_S(A_S,A_F),                                                     \\
  \alpha_{1F} & =\exp\left(\frac{h}{2}\mathcal{D}_F\right)
  \mathcal{N}_F(A_S,A_F),                                                     \\
  \alpha_{2S} & =\mathcal{N}_S\left(A_S^{\mathrm{ip}}+\frac{h}{2}\alpha_{1S},
  A_F^{\mathrm{ip}}+\frac{h}{2}\alpha_{1F}\right),                            \\
  \alpha_{2F} & =\mathcal{N}_F\left(A_S^{\mathrm{ip}}+\frac{h}{2}\alpha_{1S},
  A_F^{\mathrm{ip}}+\frac{h}{2}\alpha_{1F}\right),                            \\
  \alpha_{3S} & =\mathcal{N}_S\left(A_S^{\mathrm{ip}}+\frac{h}{2}\alpha_{2S},
  A_F^{\mathrm{ip}}+\frac{h}{2}\alpha_{2F}\right),                            \\
  \alpha_{3F} & =\mathcal{N}_F\left(A_S^{\mathrm{ip}}+\frac{h}{2}\alpha_{2S},
  A_F^{\mathrm{ip}}+\frac{h}{2}\alpha_{2F}\right),                            \\
  \alpha_{4S} & =\mathcal{N}_S\left(
  \exp\left(\frac{h}{2}\mathcal{D}_S\right)
  (A_S^{\mathrm{ip}}+h\alpha_{3S}),
  \exp\left(\frac{h}{2}\mathcal{D}_F\right)
  (A_F^{\mathrm{ip}}+h\alpha_{3F})\right),                                    \\
  \alpha_{4F} & =\mathcal{N}_F\left(
  \exp\left(\frac{h}{2}\mathcal{D}_S\right)
  (A_S^{\mathrm{ip}}+h\alpha_{3S}),
  \exp\left(\frac{h}{2}\mathcal{D}_F\right)
  (A_F^{\mathrm{ip}}+h\alpha_{3F})\right),                                    \\
  A_S(z+h)    & =\exp\left(\frac{h}{2}\mathcal{D}_S\right)
  \left[A_S^{\mathrm{ip}}+\frac{h}{6}
    (\alpha_{1S}+2\alpha_{2S}+2\alpha_{3S})\right]
  +\frac{h}{6}\alpha_{4S},                                                    \\
  A_F(z+h)    & =\exp\left(\frac{h}{2}\mathcal{D}_F\right)
  \left[A_F^{\mathrm{ip}}+\frac{h}{6}
    (\alpha_{1F}+2\alpha_{2F}+2\alpha_{3F})\right]
  +\frac{h}{6}\alpha_{4F}.
\end{align}
One should use the form of Equation~\eqref{eq:IP-CWE} and finally apply
$z'=z+h/2$ when deriving the above relations, and remember to include the
phase-matching term when doing the actual numerical calculation.

\printbibliography

\end{document}
